% === Background and Motivating Example ===
\section{Background and Motivating Example}
\label{sec:background}

Consider an engineer on call for a microservice application, paged when checkout requests begin to fail. The pressing question is not how to repair the
system from first principles but \emph{what is happening} and \emph{whether it
warrants action}. In a fully instrumented deployment the evidence is rich yet noisy: the cluster reports a message broker's pod unavailable, its consumer services begin logging errors, distributed traces show their requests failing downstream, and an availability alert fires.

The shortcut a seasoned engineer takes is \emph{recognition}: this pattern, a broker outage stalling its consumers and surfacing as checkout failures, has occurred before, and a resolved ticket already records how it was diagnosed and fixed. Retrieving that ticket turns an open-ended investigation into a near-instant analogy, the most direct path from symptom to resolution, provided the past incident can be found. We do not attempt to prove the root cause of the live incident from first principles; instead, we rank historical incidents that provide plausible analogies for human triage.

Two characteristics of operational triage shape our design. First, the \emph{same incident can arrive with vastly different evidence}. The broker failure above may carry logs, metrics, traces, and cluster events; in another team's tracker the identical failure may be a single sentence (``consumers stalled after the broker restarted, backlog growing'') with no telemetry. Yet the task is identical: identify the most relevant resolved incident. A system that depends on complete observability fails precisely when historical knowledge is needed most. Second, an incident is \emph{not instantaneous}: it evolves over a sequence of windows, from early warning signs through active failure to recovery, so treating each window independently produces repeated alerts for one evolving incident.

\smallskip
\noindent\textbf{What an incident is, concretely.} Before introducing notation we
state what the objects are, since the formalism is otherwise easy to read as
abstract. \emph{A memory item is a resolved issue-tracker ticket} written by
engineers in the course of their work. On our real corpus it is a genuine Apache
Jira issue of type \textsc{Bug} with resolution \textsc{Fixed}; the retrievable
document is its summary, description, and the log and stack-trace blocks its
authors pasted into the thread, and there are $38{,}642$ of them. On our
telemetry corpora it is a ticket generated from a deliberately injected fault's
ground truth and rewritten into engineer prose, which is why the link from
observation to ticket there is exact rather than inferred.

\emph{A query is whatever the system is handed at decision time.} In an
instrumented deployment that is a five-minute slice of one service's activity:
its log lines, its CPU, memory and restart series, any firing alerts, the
cluster's pod events, and a sample of distributed-trace spans. In a text-only
deployment it is the incident report itself, with no telemetry at all. The task
is identical in both cases, which is the point.

\smallskip
\noindent\textbf{Problem formulation.}
We formulate incident triage as adaptive retrieval over a time-aware incident memory under heterogeneous evidence. Let
$\mathcal{W}=\{w_1,\ldots,w_n\}$ denote a time-ordered stream of observation
windows, where each window
\[
w_i=(t_i,\mathbf{x}_i,e_i)
\]
is observed at time $t_i$, contains a textual summary $\mathbf{x}_i$, and may
include an arbitrary subset of auxiliary evidence
$e_i\subseteq\{\text{metrics},\text{logs},\text{traces},\text{events}\}$.

The retrieval memory
$\mathcal{M}=\{m_1,\ldots,m_N\}$ consists of previously resolved incidents.
Each incident
\[
m_j=(\mathbf{d}_j,\mathcal{E}_j,\tau_j)
\]
contains its textual description $\mathbf{d}_j$, a set of extracted entities
$\mathcal{E}_j$ (e.g., services, components, error classes, and symptoms), and
its resolution time $\tau_j$. Concretely, $\mathbf{d}_j$ is the ticket text
described above, $\mathcal{E}_j$ is produced offline by a language model reading
only that text (\S\ref{sec:approach-kg}), and $\tau_j$ is the issue's resolution
timestamp. To prevent temporal leakage, only incidents
resolved before the current observation are searchable; for a window observed at
time $t_i$, the eligible memory is
\[
\mathcal{M}_{<t_i}=\{m_j\in\mathcal{M}\mid \tau_j<t_i\}.
\]

In all evaluations, the query incident itself is removed from the searchable memory; for real Jira issues, the query time is the issue creation time and memory eligibility is determined by the candidate issue’s resolution time. Given a window $w_i$, the objective is to retrieve the top-$k$ most relevant
historical incidents from $\mathcal{M}_{<t_i}$ while simultaneously deciding
whether the window warrants an engineer page. Since consecutive windows often belong to the same operational event, the system must additionally associate related windows into one incident so that only one page is issued.

\smallskip
\noindent\textbf{Why windows.} Nothing in the formulation requires discretizing
the stream: $\mathcal{M}_{<t}$ is defined continuously in $t$, and no part of the
retrieval depends on where a window boundary falls. We window because that is the
granularity at which \emph{ground truth exists}. A label states that a competent
engineer would have opened an issue over some interval of activity, and our
synthetic corpora are annotated per five-minute interval; a continuous
formulation would need continuous supervision we cannot obtain. Windowing also
matches how the alerting rules these datasets emulate are themselves evaluated,
and it is what creates the duplicate-page problem the state layer exists to
solve. The windowing is therefore a property of the evaluation, not a commitment
of the method.
